% Options for packages loaded elsewhere
\PassOptionsToPackage{unicode}{hyperref}
\PassOptionsToPackage{hyphens}{url}
\documentclass[
]{article}
\usepackage{xcolor}
\usepackage[margin=1in]{geometry}
\usepackage{amsmath,amssymb}
\setcounter{secnumdepth}{-\maxdimen} % remove section numbering
\usepackage{iftex}
\ifPDFTeX
  \usepackage[T1]{fontenc}
  \usepackage[utf8]{inputenc}
  \usepackage{textcomp} % provide euro and other symbols
\else % if luatex or xetex
  \usepackage{unicode-math} % this also loads fontspec
  \defaultfontfeatures{Scale=MatchLowercase}
  \defaultfontfeatures[\rmfamily]{Ligatures=TeX,Scale=1}
\fi
\usepackage{lmodern}
\ifPDFTeX\else
  % xetex/luatex font selection
\fi
% Use upquote if available, for straight quotes in verbatim environments
\IfFileExists{upquote.sty}{\usepackage{upquote}}{}
\IfFileExists{microtype.sty}{% use microtype if available
  \usepackage[]{microtype}
  \UseMicrotypeSet[protrusion]{basicmath} % disable protrusion for tt fonts
}{}
\makeatletter
\@ifundefined{KOMAClassName}{% if non-KOMA class
  \IfFileExists{parskip.sty}{%
    \usepackage{parskip}
  }{% else
    \setlength{\parindent}{0pt}
    \setlength{\parskip}{6pt plus 2pt minus 1pt}}
}{% if KOMA class
  \KOMAoptions{parskip=half}}
\makeatother
\usepackage{longtable,booktabs,array}
\newcounter{none} % for unnumbered tables
\usepackage{calc} % for calculating minipage widths
% Correct order of tables after \paragraph or \subparagraph
\usepackage{etoolbox}
\makeatletter
\patchcmd\longtable{\par}{\if@noskipsec\mbox{}\fi\par}{}{}
\makeatother
% Allow footnotes in longtable head/foot
\IfFileExists{footnotehyper.sty}{\usepackage{footnotehyper}}{\usepackage{footnote}}
\makesavenoteenv{longtable}
\usepackage{graphicx}
\makeatletter
\newsavebox\pandoc@box
\newcommand*\pandocbounded[1]{% scales image to fit in text height/width
  \sbox\pandoc@box{#1}%
  \Gscale@div\@tempa{\textheight}{\dimexpr\ht\pandoc@box+\dp\pandoc@box\relax}%
  \Gscale@div\@tempb{\linewidth}{\wd\pandoc@box}%
  \ifdim\@tempb\p@<\@tempa\p@\let\@tempa\@tempb\fi% select the smaller of both
  \ifdim\@tempa\p@<\p@\scalebox{\@tempa}{\usebox\pandoc@box}%
  \else\usebox{\pandoc@box}%
  \fi%
}
% Set default figure placement to htbp
\def\fps@figure{htbp}
\makeatother
\setlength{\emergencystretch}{3em} % prevent overfull lines
\providecommand{\tightlist}{%
  \setlength{\itemsep}{0pt}\setlength{\parskip}{0pt}}
\usepackage{bookmark}
\IfFileExists{xurl.sty}{\usepackage{xurl}}{} % add URL line breaks if available
\urlstyle{same}
\hypersetup{
  hidelinks,
  pdfcreator={LaTeX via pandoc}}

\author{}
\date{\vspace{-2.5em}}

\begin{document}

\section{\texorpdfstring{Benchmarking Memo: \texttt{sndem}
vs.~\texttt{sndem\_benchmark} for Within-Country
Analysis}{Benchmarking Memo: sndem vs.~sndem\_benchmark for Within-Country Analysis}}\label{benchmarking-memo-sndem-vs.-sndem_benchmark-for-within-country-analysis}

\textbf{Date:} June 2026\\
\textbf{Prepared by:} PM, with assistance from Claude Code Sonnet 4.6\\
\textbf{For:} Michael Coppedge and co-authors

\begin{center}\rule{0.5\linewidth}{0.5pt}\end{center}

\subsection{1. What the method produces and
why}\label{what-the-method-produces-and-why}

The snvdem index is built from two sources of information that serve
different roles:

\textbf{Expert data (weights):} Four V-Dem survey questions ask country
experts what characteristics are associated with (1) more free and fair
elections, (2) less free and fair elections, (3) stronger civil
liberties, and (4) weaker civil liberties. The proportions of coder
agreement around each characteristic and question are synthesized into
relevance weights for each dimension --- elections (\texttt{snelect})
and civil liberties (\texttt{sncivlib}).

\textbf{Objective data (geopredictors):} Geocoded structural variables
measured at the municipal level in Colombia (rurality, distance from
capital, population density, violence indicators, etc.) are standardized
using an empirical CDF applied across the full 2000--2023 panel. This
maps each predictor to {[}0,1{]} representing its percentile within the
full Colombian distribution.

\textbf{The intermediate index:} Expert weights are applied to the
CDF-standardized geopredictors, producing a weighted average for each
dimension. The paper's Equation 1 then defines:

\begin{verbatim}
sndem = 0.5 * (snelect + sncivlib)
\end{verbatim}

Both subindices naturally fall on {[}0,1{]} because the inputs are
CDF-standardized. The benchmark step (folder 08) is then applied to each
dimension separately.

\begin{center}\rule{0.5\linewidth}{0.5pt}\end{center}

\subsection{2. A critical reading of the operational strategy
document}\label{a-critical-reading-of-the-operational-strategy-document}

Two excerpts from the revised operational strategy (Jan 2026) are in
tension with each other, and that tension is creating confusion around
our data and benchmarking.

\subsubsection{Excerpt 1 --- a conceptual definition
(p.~1)}\label{excerpt-1-a-conceptual-definition-p.-1}

\begin{quote}
``We conceive of free and fair elections at the municipal level as
\textbf{a set of deviations centered around this summary subnational
measure} {[}v2elffelr{]}. In order to benchmark the many municipal-level
values, we need to determine (a) the maximum range of variation around
the summary measure in the entire country in a given year and (b) how
high or low each municipality is within that range.''
\end{quote}

This is making a \textbf{theoretical claim} about what municipal
election quality \emph{is}: a deviation from the national V-Dem
subnational average. If taken at face value, this means
\texttt{EL\_col\_mt} --- each municipality's position expressed as a
deviation from \texttt{v2elffelr} within the expert-estimated national
range --- is the conceptually intended final measure, not an optional
calibration.

This matters for temporal interpretation: two municipalities with the
same unbenchmarked \texttt{snelect\ =\ 0.6} in different years are not
equivalent if the national V-Dem anchor (\texttt{v2elffelr}) changed
between those years. The benchmark anchors each year's scores to the
current national estimate, making scores comparable over time within
Colombia as well as across countries.

\subsubsection{Excerpt 2 --- a stated purpose
(p.~3)}\label{excerpt-2-a-stated-purpose-p.-3}

\begin{quote}
``To complete the calibration of the data \textbf{for cross-national
comparisons}, the subnational estimates such as those we have generated
for Colombia can be centered around the means of their respective series
and rescaled to fit within their ranges.''
\end{quote}

This phrase appears to restrict the purpose of benchmarking to
cross-national comparison only. But in context, ``for cross-national
comparisons'' may be describing an additional payoff rather than the
exclusive purpose. The preceding sentence reads: ``For \emph{both} civil
liberties and free and fair elections, the characteristics of each
location determine how much it varies within that maximum national
range'' --- a general statement about the measure, not a
cross-national-only claim.

\subsubsection{The tension}\label{the-tension}

{\def\LTcaptype{none} % do not increment counter
\begin{longtable}[]{@{}
  >{\raggedright\arraybackslash}p{(\linewidth - 2\tabcolsep) * \real{0.5000}}
  >{\raggedright\arraybackslash}p{(\linewidth - 2\tabcolsep) * \real{0.5000}}@{}}
\toprule\noalign{}
\begin{minipage}[b]{\linewidth}\raggedright
Excerpt
\end{minipage} & \begin{minipage}[b]{\linewidth}\raggedright
Implication
\end{minipage} \\
\midrule\noalign{}
\endhead
\bottomrule\noalign{}
\endlastfoot
``We conceive of free and fair elections as deviations centered around
v2elffelr'' & Benchmarked scores ARE the intended final measure --- for
any analytical use \\
``To complete the calibration for cross-national comparisons'' &
Benchmarking is supplementary, needed only for cross-national
comparison \\
\end{longtable}
}

These two readings lead to different practical conclusions about which
variable to use for within-country mapping and regression.

\begin{center}\rule{0.5\linewidth}{0.5pt}\end{center}

\subsection{3. The unsolved combination
problem}\label{the-unsolved-combination-problem}

Both excerpts describe the benchmark in terms of \textbf{individual
dimensions}: \texttt{EL\_col\_mt} (elections on the V-Dem latent scale,
\textasciitilde{[}-3.5, 3.5{]}) and \texttt{CL\_col\_mt} (civil
liberties on {[}0,1{]}, anchored to \texttt{v2x\_civlib}). The
operational strategy document shows these as separate figures (Figures 4
and 5). It does not describe how to combine them into a single
\texttt{sndem\_benchmark}.

This is the implementation gap. Every approach Patrick tried in June
2026 produced distortions:

{\def\LTcaptype{none} % do not increment counter
\begin{longtable}[]{@{}
  >{\raggedright\arraybackslash}p{(\linewidth - 2\tabcolsep) * \real{0.5000}}
  >{\raggedright\arraybackslash}p{(\linewidth - 2\tabcolsep) * \real{0.5000}}@{}}
\toprule\noalign{}
\begin{minipage}[b]{\linewidth}\raggedright
Approach
\end{minipage} & \begin{minipage}[b]{\linewidth}\raggedright
Problem
\end{minipage} \\
\midrule\noalign{}
\endhead
\bottomrule\noalign{}
\endlastfoot
Within-year \texttt{ecdf()} (original script) & V-Dem anchor
mathematically cancels; national average fixed at 0.5 every year;
temporal trend removed \\
Within-bounds normalization & Moving reference frame: V-Dem anchors rise
with Colombia's improvement; CDF-based geopredictors do not; inverted
trend \\
\texttt{pnorm(EL\_col\_mt)} & Trend direction preserved but non-linearly
amplified (pnorm is steepest near 0, where Colombia's EL scores
cluster) \\
Cross-panel linear normalization & Still exaggerated because EL and CL
capture different magnitudes of absolute change \\
\end{longtable}
}

\textbf{Root cause:} The geopredictors' CDF is computed over the full
2000--2023 panel, capturing cross-sectional relative position but not
absolute temporal improvement. The V-Dem anchors (\texttt{v2elffelr},
\texttt{CLSNmean}) do track absolute improvement over time. Every
normalization that tries to merge these two types of information either
removes the temporal trend or distorts its shape.

\begin{center}\rule{0.5\linewidth}{0.5pt}\end{center}

\subsection{4. A structural note on the geopredictors and
time}\label{a-structural-note-on-the-geopredictors-and-time}

Because the CDF is applied across the full panel (all municipality-years
from 2000--2023), each predictor represents a municipality's
\emph{percentile rank within the full Colombian distribution}, not its
absolute level. This design choice enables replication across countries
with very different absolute levels --- but it means \texttt{sndem} is
primarily a cross-sectional measure. Apparent trends over time reflect
changes in which municipalities rank higher or lower on the weighted
predictors, not absolute democratic improvement.

This is arguably acceptable for within-country regression (which asks:
``what predicts relative democratic quality across municipalities?'')
but creates difficulty when combined with a time-varying V-Dem anchor in
the benchmark step. The benchmark is designed to capture absolute level,
while the geopredictors capture relative rank --- these are not the same
thing.

\begin{center}\rule{0.5\linewidth}{0.5pt}\end{center}

\subsection{5. Core questions requiring
clarification}\label{core-questions-requiring-clarification}

The following questions cannot be resolved by implementation choices
alone. They require a methodological decision by the authors:

\textbf{Q1 --- Is Excerpt 1 a theoretical commitment or a motivating
framework?}

Does ``we conceive of free and fair elections as deviations centered
around v2elffelr'' mean that \texttt{EL\_col\_mt} (the benchmarked
score) is the \emph{intended final measure} for all uses including
within-country mapping and regression? Or is it describing the
theoretical logic that motivates a calibration step whose outputs
(\texttt{EL\_col\_mt}, \texttt{CL\_col\_mt}) are only needed when
comparing across countries?

\emph{This determines whether prior descriptive maps and regressions
using \texttt{sndem} are using the correct variable.}

\textbf{Q2 --- How should the two benchmarked dimensions be combined?}

The methodology documents describe \texttt{EL\_col\_mt} and
\texttt{CL\_col\_mt} as outputs per dimension (separate figures) but do
not specify how to combine them into a single \texttt{sndem\_benchmark}.
They are on incompatible scales by design (latent elections scale
vs.~{[}0,1{]} civil liberties scale). Any normalization step introduces
either scale distortion or trend distortion. Is there a principled
combination rule the authors intend?

\emph{This is a calculation bottleneck.}

\textbf{Q3 --- Does the CDF standardization limit temporal
interpretation?}

The full-panel CDF makes \texttt{sndem} a cross-sectional measure. Is it
the intent of the methodology that the index captures relative
democratic quality at each point in time (cross-sectional validity),
with temporal trends being a secondary and potentially noisy feature? Or
should a different standardization approach be used to better preserve
temporal information?

\emph{This affects how time-series maps and trend figures in the paper
should be framed.}

\begin{center}\rule{0.5\linewidth}{0.5pt}\end{center}

\subsection{6. Practical recommendations pending
clarification}\label{practical-recommendations-pending-clarification}

\textbf{If Excerpt 1 is a theoretical commitment (benchmarking is
required for all uses):}\\
The combination problem (Q2) must be resolved before any maps or
regressions proceed. The current \texttt{01\_benchmark.R} outputs
(\texttt{EL\_col\_mt}, \texttt{CL\_col\_mt}) are the right intermediate
products; a principled combination rule is needed.

\textbf{If Excerpt 1 is a motivating framework (benchmarking is for
cross-national use only):}\\
Use \texttt{sndem\ =\ 0.5\ *\ (snelect\ +\ sncivlib)} from the weighting
script for all within-country analysis. Prior descriptive maps and
regressions using this variable are methodologically sound.
\texttt{EL\_col\_mt} and \texttt{CL\_col\_mt} are supplementary outputs
for cross-national comparison.

\textbf{Either way}, the paper's current language should be made
internally consistent: the theoretical framing in Excerpt 1 and the
stated purpose in Excerpt 2 should not point in different directions.

\end{document}
